
\documentclass{article}
\usepackage{colortbl}
\usepackage{makecell}
\usepackage{multirow}
\usepackage{supertabular}

\begin{document}

\newcounter{utterance}

\centering \large Interaction Transcript for game `multimodal\_referencegame', experiment `CLEVR\_static\_target\_images', episode 1 with NEW\_Reflect\_Agent.
\vspace{24pt}

{ \footnotesize  \setcounter{utterance}{1}
\setlength{\tabcolsep}{0pt}
\begin{supertabular}{c@{$\;$}|p{.15\linewidth}@{}p{.15\linewidth}p{.15\linewidth}p{.15\linewidth}p{.15\linewidth}p{.15\linewidth}}
    \# & $\;$A & \multicolumn{4}{c}{Game Master} & $\;\:$B\\
    \hline

    \theutterance \stepcounter{utterance}  
    & & \multicolumn{4}{p{0.6\linewidth}}{
        \cellcolor[rgb]{0.9,0.9,0.9}{
            \makecell[{{p{\linewidth}}}]{
                \texttt{\tiny{[P1$\langle$GM]}}
                \texttt{You are given three images, one is called target and the other two are distractors.} \\
\texttt{Your task is to generate a referring expression that best describes the target image while distinguishing it from the two other distractor images.} \\
\texttt{The first image is a distractor, the second image is a distractor, and the third image is the target. Instruction: Describe the target image. Generate the referring expression starting with the tag "Expression: " for the given target image. Omit any other text.} \\
            }
        }
    }
    & \\ \\

    \theutterance \stepcounter{utterance}  
    & \multicolumn{4}{p{0.45\linewidth}}{
        \cellcolor[rgb]{0.95,0.95,0.95}{
            \makecell[{{p{\linewidth}}}]{
                \texttt{\tiny{[P1$|$internal]}}
                \texttt{Reflections from previous episodes:} \\
\texttt{{-}} \\
\texttt{ Failed. I generated a generic description ("a close{-}up of a red apple") that didn't distinguish the target from the distractors (which likely also showed red apples). I didn't analyze the key *differences* between the images. Next time, I will systematically compare all three images first, identify the target's unique feature (e.g., a specific stem, leaf, or bite mark), and explicitly include that distinguishing detail in the expression.} \\
\texttt{{-}} \\
\texttt{ Failed. I did not generate an expression starting with "Expression: " as instructed. I omitted the required tag. Next time, I will strictly follow the output format by beginning my response with "Expression: " followed by the description.} \\
\texttt{{-}} \\
\texttt{ Failed. I generated a generic description ("a person sitting at a table with a laptop") that didn't distinguish the target from the distractors (which also showed people with laptops). I didn't analyze the key *distinctive* feature (e.g., the target's specific clothing, pose, or background). Next time, I will systematically compare all three images first, identify the unique attribute(s) of the target, and explicitly include that contrast in the expression.} \\
            }
        }
    }
    & & \\ \\

    \theutterance \stepcounter{utterance}  
    & \multicolumn{4}{p{0.45\linewidth}}{
        \cellcolor[rgb]{0.95,0.95,0.95}{
            \makecell[{{p{\linewidth}}}]{
                \texttt{\tiny{[P1$|$internal]}}
                \texttt{{[}attempt 1 – valid{]}} \\
\texttt{Expression: a large blue sphere in the foreground with a small glowing blue cylinder next to it, distinct from the other images which have different arrangements of objects and lack this specific combination.} \\
            }
        }
    }
    & & \\ \\

    \theutterance \stepcounter{utterance}  
    & \multicolumn{4}{p{0.6\linewidth}}{
        \cellcolor[rgb]{0.9,0.9,0.9}{
            \makecell[{{p{\linewidth}}}]{
                \texttt{\tiny{[P1$\rangle$GM]}}
                \texttt{Expression: a large blue sphere in the foreground with a small glowing blue cylinder next to it, distinct from the other images which have different arrangements of objects and lack this specific combination.} \\
            }
        }
    }
    & & \\ \\

    \theutterance \stepcounter{utterance}  
    & & & \multicolumn{2}{p{0.3\linewidth}}{
        \cellcolor[rgb]{0.9,0.9,0.9}{
            \makecell[{{p{\linewidth}}}]{
                \texttt{\tiny{[GM$|$GM]}}
                \texttt{Expression: a large blue sphere in the foreground with a small glowing blue cylinder next to it, distinct from the other images which have different arrangements of objects and lack this specific combination.} \\
            }
        }
    }
    & & \\ \\

    \theutterance \stepcounter{utterance}  
    & & \multicolumn{4}{p{0.6\linewidth}}{
        \cellcolor[rgb]{0.9,0.9,0.9}{
            \makecell[{{p{\linewidth}}}]{
                \texttt{\tiny{[GM$\rangle$P2]}}
                \texttt{You are given three images.} \\
\texttt{You are also given a referring expression that describes one of the given images.} \\
\texttt{Your task is to select the image that matches the given referring expression.} \\
\texttt{Generate only the number (in text) of the image that matches the given expression by selecting first, second, or third.} \\
\\ 
\texttt{Expression: a large blue sphere in the foreground with a small glowing blue cylinder next to it, distinct from the other images which have different arrangements of objects and lack this specific combination.} \\
\texttt{Question: Which image does the expression refer to?} \\
\texttt{Start with the tag "Answer: ", followed by your selection. Omit any other text.} \\
            }
        }
    }
    & \\ \\

    \theutterance \stepcounter{utterance}  
    & & & \multicolumn{4}{p{0.45\linewidth}}{
        \cellcolor[rgb]{0.95,0.95,0.95}{
            \makecell[{{p{\linewidth}}}]{
                \texttt{\tiny{[P2$|$internal]}}
                \texttt{Reflections from previous episodes:} \\
\texttt{{-}} \\
\texttt{ Answer: second} \\
\\ 
\texttt{Reflection: Succeeded. The expression clearly described a large blue cube as the main object, with specific smaller objects in the background. I matched these attributes to the correct image. Next time, I will continue to carefully parse spatial relationships ("in the background") and color/shape combinations to ensure accurate selection.} \\
\texttt{{-}} \\
\texttt{ Failed. The expression described a specific spatial arrangement ("next to") and background objects that likely didn't match the actual image layout. I guessed without being able to verify the visual details. Next time, I would need visual input to compare the description to the actual images.} \\
\texttt{{-}} \\
\texttt{ Answer: second} \\
\\ 
\texttt{Reflection: Succeeded. The expression specified a large red metallic cylinder with a small purple cylinder nearby. Image 2 matched both objects and their spatial relationship. I correctly identified the key attributes (size, color, material, proximity) and ignored distractors. Next time, I will continue to systematically verify each descriptive detail against the images.} \\
            }
        }
    }
     \\ \\

    \theutterance \stepcounter{utterance}  
    & & & \multicolumn{4}{p{0.45\linewidth}}{
        \cellcolor[rgb]{0.95,0.95,0.95}{
            \makecell[{{p{\linewidth}}}]{
                \texttt{\tiny{[P2$|$internal]}}
                \texttt{{[}attempt 1 – valid{]}} \\
\texttt{Answer: second} \\
            }
        }
    }
     \\ \\

    \theutterance \stepcounter{utterance}  
    & & & \multicolumn{4}{p{0.6\linewidth}}{
        \cellcolor[rgb]{0.9,0.9,0.9}{
            \makecell[{{p{\linewidth}}}]{
                \texttt{\tiny{[GM$\langle$P2]}}
                \texttt{Answer: second} \\
            }
        }
    }
     \\ \\

    \theutterance \stepcounter{utterance}  
    & & & \multicolumn{2}{p{0.3\linewidth}}{
        \cellcolor[rgb]{0.9,0.9,0.9}{
            \makecell[{{p{\linewidth}}}]{
                \texttt{\tiny{[GM$|$GM]}}
                \texttt{second} \\
            }
        }
    }
    & & \\ \\

\end{supertabular}
}

\end{document}
